% Kapitel: Das Automotive-Lighting-Plugin (Scheinwerfer und Signalleuchten mit Glühlampen)
% Quellen: python/plugins/automotive_lighting/ (al_*.py, plugin.json, README.md), CHANGELOG.md
\chapter{Scheinwerfer und Leuchten: Automotive Lighting}\label{kap:automotive}
\index{Automotive Lighting}\index{Scheinwerfer}

Das Plugin \plugin{Automotive Lighting} entwirft Scheinwerfer und Signalleuchten für Fahrzeuge,
deren Lichtquelle eine Glühlampe ist -- eine H7, eine H4, eine P21W und viele andere. Es erzeugt
die optischen Flächen als exakte NURBS-Flächen im Modell, berechnet die Wendelbilder und die
Lichtverteilung auf dem Messschirm, prüft sie gegen die Prüfpunkte der Regelungen UN/ECE R112
(Abblendlicht) und UN R148 (Signalleuchten) und kontrolliert den Entwurf mit einem Raytracing, das
die Lampe so nachbildet, wie sie ist: Wendel als Drahtwendel, Glaskolben, schwarze Kappe, der Schirm
der H4.

Das Plugin kennt zwei Bauarten (\gui{System} im Panel):

\begin{description}
  \item[Klarglasscheinwerfer mit facettiertem Reflektor] Der Spiegel selbst formt das Licht und
    seine \textbf{Hell-Dunkel-Grenze}. Er ist in Zonen und Streifen (Facetten) geteilt; jede Facette
    legt die Wendelbilder ihrer Punkte genau an die Hell-Dunkel-Grenze. Wände mit Entformungsschräge
    und Verrundungen verbinden die Facetten zu einer Fläche, die sich spritzgießen lässt.
  \item[Scheinwerfer mit Streuscheibe] Ein Paraboloidspiegel und davor eine Streuscheibe, deren
    Innenseite ein Raster von Elementen trägt -- Prismen, Zylinderlinsen, Kissen --, die das Licht
    lenken und verbreitern. Die Hell-Dunkel-Grenze kommt hier vom Schirm der Lampe (H4); das Plugin
    kann die Elemente der Scheibe automatisch optimieren.
\end{description}

Dieses Kapitel führt durch alle Funktionen. Die vollständigen Parametertabellen, die Lampendaten,
die Prüfpunkte und alle Ergebnisse stehen im Anhang, \autoref{kap:plugin-automotive}.

\section{Überblick}\label{sec:al-ueberblick}

\subsection{Voraussetzungen}\label{sec:al-voraussetzungen}

Das Plugin (Kennung \code{al}, Name \plugin{Automotive Lighting}) wird mit der Anwendung
installiert. Für Entwurf, Wendelbilder und Lichtverteilung braucht es nur \code{numpy} in der
Python-Umgebung der Plugins. Das Raytracing (\gui{Light distribution (ray tracing)}, die Unschärfe
bei \gui{Optimise lens}) verwendet den Rechenkern des Plugins \plugin{Path Tracer}
(\autoref{kap:pathtracer}):

\begin{itemize}
  \item Ist das Paket \code{pathtracer} in Version 0.3.0 oder neuer installiert, rechnet der
    C++-Kern -- etwa fünf Millionen Lichtwege durch Kolben und facettierten Reflektor in rund
    zehn Sekunden.
  \item Sonst rechnet der NumPy-Kern aus dem Ordner des Path-Tracer-Plugins, der neben diesem Plugin
    liegen muss. Er braucht für einige hunderttausend Lichtwege Minuten.
\end{itemize}

\gui{Create environment…} in \menu{Plugins,Settings…} installiert beides (siehe
\autoref{kap:installation} und \autoref{kap:plugins}).

\subsection{Die Aktionen}

\begin{referenz}{>{\raggedright\arraybackslash}p{4.4cm}>{\raggedright\arraybackslash}X}
\toprule
Aktion & Wirkung \\
\midrule
\endhead
\gui{Generate reflector} & erzeugt Spiegel (Facetten und Verbindungsflächen bzw. Paraboloid),
  Streuscheibe und Wendelzylinder als NURBS-Flächen im Modell \\
\gui{Filament images} & zeigt die Umrisse der Wendelbilder vieler Spiegelpunkte auf dem Messschirm,
  mit der Hell-Dunkel-Grenze \\
\gui{Light distribution} & Lichtstärkeverteilung, Tabelle der Prüfpunkte, zwei Schnitte, Lichtströme
  und Schärfe der Hell-Dunkel-Grenze \\
\gui{Light distribution (ray tracing)} & dasselbe durch Raytracing mit der Lampe, wie sie ist \\
\gui{Generate lamp} & legt die Lampe des Raytracings (Drahtwendel, Kolben, Kappe, H4-Schirm) als
  NURBS-Flächen ins Modell \\
\gui{Optimise lens} & wählt für jede Zelle der Streuscheibe ihr Element so, dass das Abblendlicht
  die Prüfpunkte mit Reserve einhält \\
\gui{Show strips} & zeigt die Startkurven und einige charakteristische Streifen jeder Facette als
  Linien \\
\gui{Default zones} & füllt die Zonentabelle für die gewählte Verkehrsseite \\
\bottomrule
\end{referenz}

Alle Aktionen finden Sie unter \menu{Plugins,Automotive Lighting}, als Knöpfe unten im Dock
\gui{Plugin} (nachdem Sie oben \gui{Automotive Lighting} gewählt haben) und in der Gruppe
\gui{Automotive Lighting} der Werkzeugspalte. Eine laufende Aktion brechen Sie mit \gui{Cancel} in
der Statusleiste ab. Alle Änderungen am Modell, die eine Aktion vornimmt, sind ein einziger
Rückgängig-Schritt.

\subsection{Wo die Ergebnisse erscheinen}

\begin{itemize}
  \item \textbf{Flächen} im Modell, in einer eigenen Gruppe des Modellbaums (Dock \gui{Model}):
    \enquote{Reflector <Lampe>}, \enquote{Headlamp <Lampe>} oder \enquote{Lamp <Lampe>}.
  \item \textbf{Bilder} (Lichtstärkeverteilung, Wendelbilder) im Fenster \gui{Image}, mit H und V in
    Grad als Achsen.
  \item \textbf{Tabellen} (Prüfpunkte) im Fenster \gui{Table}.
  \item \textbf{Kurven} (Schnitte durch die Verteilung) im Fenster \gui{Plot}.
  \item \textbf{Linien} (Streifen) als Ergebnis in den Ansichten.
  \item \textbf{Protokoll} im Dock \gui{Plugin log}: Flächen, Lichtströme, Schärfe, Warnungen.
  \item eine \textbf{Zusammenfassung} in der Statusleiste.
\end{itemize}

Jede Rechnung legt im Dock \gui{Results} einen Ergebnissatz an, etwa \enquote{Low beam H7} oder
\enquote{Low beam H7 – ray tracing}; so können Sie verschiedene Stände nebeneinander vergleichen.

\section{Koordinaten und Messschirm}\label{sec:al-koordinaten}
\index{Messschirm}

\paragraph{Im Modell.} Das Plugin rechnet in Millimetern mit
\begin{itemize}
  \item $+X$ in Fahrtrichtung (das Licht verlässt den Scheinwerfer in $+X$),
  \item $+Y$ nach links,
  \item $+Z$ nach oben.
\end{itemize}
Der Mittelpunkt der Wendel liegt im Punkt \gui{Filament centre} (Gruppe \gui{Lamp}, Standard
Ursprung). Alle erzeugten Flächen werden relativ zu diesem Punkt gelegt.

\paragraph{Auf dem Messschirm.} Die Lichtverteilung wird über zwei Winkel dargestellt: H
(horizontal, positiv nach \textbf{rechts}) und V (vertikal, positiv nach \textbf{oben}), beide in
Grad. Die Lichtstärke ist in Candela (cd) angegeben; 625 cd entsprechen einer Beleuchtungsstärke von
1 lx in 25 m Abstand. Die Linie H = 0 ist die Fahrzeugachse, V = 0 der Horizont.

\paragraph{Azimut.} Zonen und Sektoren werden über einen Azimut um die Lampenachse beschrieben, in
Fahrtrichtung gesehen: \ang{0} = links, \ang{90} = oben, \ang{180} = rechts, \ang{270} = unten.

\section{Das Panel}\label{sec:al-panel}

Das Panel ist in Gruppen gegliedert, in dieser Reihenfolge:

\begin{referenz}{>{\raggedright\arraybackslash}p{3.2cm}>{\raggedright\arraybackslash}X}
\toprule
Gruppe & Inhalt \\
\midrule
\endhead
\gui{System} & Bauart: Klarglas mit facettiertem Reflektor oder Paraboloid mit Streuscheibe \\
\gui{Lamp} & Lampentyp, eigene Lampendaten, Lage der Wendel \\
\gui{Cut-off} & Verkehrsseite, Höhe, Anstieg und Knick der Hell-Dunkel-Grenze \\
\gui{Frame} & Umriss des Spiegels (Kurven oder gerundetes Rechteck), Lampenloch \\
\gui{Reflector} & Brennweite, Streifenbreite, Facettenlänge, Zonen, Knotenabstand \\
\gui{Lens} & Lage, Form und Material der Streuscheibe, Zellen, Sektoren, Elemente, Einstellungen von
  \gui{Optimise lens} (nur Streuscheibe) \\
\gui{Connecting faces} & Verbindungsflächen an/aus, Entformungsschräge, Verrundungsradius \\
\gui{Analysis} & Reflexionsgrade, Prüfpunkte, Schatten, Direktlicht, Abtastung, Zahl der
  Wendelbilder \\
\gui{Ray tracing} & Lichtwege, Weglänge, Wendel, Kolben, H4-Schirm, Rechenkern \\
\bottomrule
\end{referenz}

Die Einstellungen werden mit dem Dokument gespeichert. \gui{Reset parameters} unten im Panel setzt
alle auf die Standardwerte zurück -- auch die Zonen- und Elementtabellen.

\begin{tipp}
Das Plugin merkt sich den zuletzt berechneten Reflektor bzw. die zuletzt berechnete Streuscheibe.
Rufen Sie nach \gui{Generate reflector} \gui{Light distribution} auf, ohne etwas an Lampe,
Hell-Dunkel-Grenze, Rahmen, Reflektor, Streuscheibe oder Verbindungsflächen zu ändern, wird der
Entwurf nicht neu gerechnet. Einstellungen der Gruppen \gui{Analysis} und \gui{Ray tracing}
erzwingen keinen neuen Entwurf.
\end{tipp}

% ---------------------------------------------------------------------------------------------
\section{Die Lampe}\label{sec:al-lampe}
\index{Glühlampe}\index{UN/ECE R37}\index{Wendel}

\subsection{Lampentypen}

Für den Entwurf ist die Wendel ein gleichmäßig leuchtender Zylinder mit Länge, Durchmesser und
Lichtstrom. \gui{Lamp} bietet diese Typen an:

\begin{center}
\small
\begin{tabular}{lrrrrl}
\toprule
Lampe & Länge & Durchmesser & Lichtstrom & Leistung & Wendel \\
 & mm & mm & lm & W & \\
\midrule
H1 & 5,0 & 1,4 & 1550 & 55 & axial \\
H4 (low beam) & 5,5 & 1,5 & 1000 & 55 & axial, über dem Schirm \\
H7 & 4,1 & 1,3 & 1500 & 55 & axial \\
H8 & 3,7 & 1,2 & 800 & 35 & axial \\
H9 & 4,8 & 1,4 & 2100 & 65 & axial \\
H11 & 4,5 & 1,3 & 1350 & 55 & axial \\
HB3 & 5,1 & 1,4 & 1860 & 60 & axial \\
HB4 & 5,1 & 1,4 & 1095 & 51 & axial \\
H3 & 5,0 & 1,4 & 1450 & 55 & quer (Nebelscheinwerfer) \\
P21W & 6,0 & 1,0 & 460 & 21 & quer \\
PY21W & 6,0 & 1,0 & 280 & 21 & quer, gelb \\
R10W & 4,0 & 0,8 & 125 & 10 & quer \\
R5W & 4,0 & 0,7 & 50 & 5 & quer \\
\bottomrule
\end{tabular}
\end{center}

Lichtstrom und Wendellänge sind die Nennwerte nach UN/ECE R37 (Rev.~7) für 12-V-Lampen: der
Lichtstrom als Sollwert bei der Prüfspannung (13,2 V, Signallampen 13,5 V), die Länge die Nennlänge
des Lampenblatts (R5W, R10W: dort keine, typische Werte). Die Durchmesser gibt R37 nicht an -- es
sind typische Werte. Standard ist die \textbf{H7}.

\begin{warnung}
Prüfen Sie die Wendelmaße gegen das Datenblatt Ihrer Lampe. Weichen sie ab, wählen Sie
\gui{Custom} und tragen die eigenen Werte ein.
\end{warnung}

\paragraph{Axial oder quer.} Bei den meisten Scheinwerferlampen liegt die Wendel längs der
Lampenachse (X). Bei H3, P21W, PY21W, R10W und R5W liegt sie quer dazu, waagerecht (entlang Y).

\subsection{Eigene Lampe}

Mit \gui{Lamp} = \gui{Custom} gelten die Felder
\begin{itemize}
  \item \gui{Custom: filament length} (Standard 4,1 mm),
  \item \gui{Custom: filament diameter} (1,3 mm),
  \item \gui{Custom: luminous flux} (1500 lm),
  \item \gui{Custom: power} (55 W bei 12 V; nur für die Drahtwendel des Raytracings),
  \item \gui{Custom: filament across the axis} (aus; an: Wendel quer, waagerecht).
\end{itemize}

\subsection{Lage der Wendel}

\gui{Filament centre} ist der Mittelpunkt der Wendel im Modell. Mit \gui{Pick} wählen Sie ihn per
Klick in einer Ansicht (mit Objektfang; \keys{Esc} bricht ab). Alle Flächen des Plugins werden
relativ zu diesem Punkt erzeugt; die Lampenachse ist immer parallel zu $+X$.

\subsection{Die H4: Schirm und Kappe}\label{sec:al-h4}
\index{H4}

Die H4 hat zwei Wendeln: die Abblendwendel über einem Schirm und dahinter die Fernlichtwendel. Der
Schirm bestimmt die Hell-Dunkel-Grenze des Abblendlichts. Das Plugin bildet ihn nach den
Lampenblättern H4/1 bis H4/5 von UN/ECE R37 nach:

\begin{itemize}
  \item Die Kanten des Schirms liegen auf der Ebene H und \ang{15} darunter; das Licht verlässt den
    Schirm über diese Kanten in die oberen \ang{195}.
  \item Die Oberkante der Wendel liegt 0,6 mm über der Ebene H.
  \item Der Schirm ist eine Wanne um die Bezugsachse, von 1,2 mm hinter der Wendel (dort durch eine
    Stirnwand geschlossen) bis 3,4 mm vor ihr (offen), Radius 4,3 mm (etwa 8,6 mm breit und 4,2 mm
    tief).
  \item Die schwarze Kappe auf der Kolbenspitze erscheint vom Achspunkt $p = \qty{28.95}{mm}$
    unter \ang{40} gegen die Achse (der Grenzwert des Blatts), auf einem Kolben von 13 mm
    Durchmesser.
\end{itemize}

Die Regelung legt Kanten, Wendellage und Kappenwinkel fest. Radius und Länge der Wanne stammen aus
der Zeichnung des Blatts H4/3, die für den Schirm nicht verbindlich ist; der Kolbendurchmesser ist
angenommen.

\begin{tipp}
In ihrem Scheinwerfer hat der Spiegel einer H4 seinen Brennpunkt an der Fernlichtwendel,
\qty{6.7}{mm} hinter der Mitte der Abblendwendel. Für einen Streuscheiben-Scheinwerfer mit H4 setzen
Sie deshalb \gui{Mirror focus behind the filament centre} = 6.7 (\autoref{sec:al-streuscheibe}).
Dann bleibt auch das Licht, das der glänzende Schirm zurückwirft, unter der Hell-Dunkel-Grenze.
\end{tipp}

% ---------------------------------------------------------------------------------------------
\section{Die Hell-Dunkel-Grenze}\label{sec:al-cutoff}
\index{Hell-Dunkel-Grenze}

Das asymmetrische Abblendlicht hat eine Hell-Dunkel-Grenze aus zwei Ästen:

\begin{itemize}
  \item einem \textbf{waagerechten Ast} auf der Seite des Gegenverkehrs, in der Höhe
    \gui{Horizontal part at V} (Standard $-\ang{0.57}$, das ist 1\,\% Neigung unter den Horizont),
  \item einem \textbf{ansteigenden Ast} auf der eigenen Seite, der vom Knick bei \gui{Kink at H}
    (Standard \ang{0}) unter dem Winkel \gui{Rise} (Standard \ang{15}) ansteigt.
\end{itemize}

\gui{Traffic} wählt die Verkehrsseite: bei \code{right-hand} (Rechtsverkehr, Standard) steigt die
Grenze nach rechts an, bei \code{left-hand} nach links. Die Prüfpunkte werden für Linksverkehr an
V-V gespiegelt.

Jeder Ast ist auf dem Messschirm eine Gerade, im Raum also eine Ebene durch die Lampe. Die
Facetten des Klarglasreflektors legen ihre Wendelbilder an einen der beiden Äste
(\autoref{sec:al-zonen}).

\begin{tipp}
Für einen Nebelscheinwerfer mit waagerechter Grenze setzen Sie \gui{Rise} = 0.
\end{tipp}

% ---------------------------------------------------------------------------------------------
\section{Der Rahmen}\label{sec:al-rahmen}

Der Rahmen ist der Umriss des Spiegels in der Vorderansicht (entlang X gesehen). Er begrenzt
Reflektor und Streuscheibe.

\begin{description}
  \item[\gui{Frame (curves)}] Geschlossene Kurven, die Sie im CAD gezeichnet haben: wählen Sie sie
    aus und drücken Sie \gui{Use selection}. Die Kurven werden Ende an Ende zu einem Umriss verbunden;
    ihre Lage wird relativ zu \gui{Filament centre} in der Y-Z-Ebene ausgewertet (die X-Koordinate
    spielt keine Rolle).
  \item[gerundetes Rechteck] Ist \gui{Frame (curves)} leer, gilt ein Rechteck mit \gui{Width}
    (Standard 170 mm), \gui{Height} (100 mm) und \gui{Corner radius} (20 mm), dessen Mitte um
    \gui{Centre to the left of the lamp axis} und \gui{Centre above the lamp axis} (beide 0) von der
    Lampenachse verschoben ist.
  \item[\gui{Lamp hole radius}] Radius des Lochs für die Lampe in der Spiegelmitte (Standard 12 mm).
\end{description}

\begin{hinweis}
Der Rahmen sollte konvex oder nahezu konvex sein. Kurven, die sich nicht zu einem Umriss schließen,
weist das Plugin mit der Lücke in Millimetern ab (\enquote{the frame curves do not join to one outline
(gap of … mm) – select curves that meet end to end}).
\end{hinweis}

% ---------------------------------------------------------------------------------------------
\section{Der Klarglasreflektor}\label{sec:al-klarglas}
\index{Klarglasreflektor}\index{facettierter Reflektor}

Mit \gui{System} = \gui{clear lens: facetted reflector} (Standard) entwirft das Plugin einen
facettierten Spiegel für einen Scheinwerfer mit klarer Abschlussscheibe. Der Spiegel allein erzeugt
die Lichtverteilung \emph{und} die Hell-Dunkel-Grenze.

\subsection{Das Prinzip: Wendelbilder an der Grenze}\label{sec:al-prinzip}
\index{Wendelbild}

Ein Spiegelpunkt $P$ mit der Normalen $n$ wirft das Licht jedes Wendelpunkts $Q$ in eine Richtung
auf den Messschirm. Alle diese Richtungen zusammen sind das \textbf{Wendelbild} des Spiegelpunkts --
ein kleiner, länglicher Fleck. Damit das Abblendlicht niemanden blendet, muss jedes Wendelbild
\emph{unter} der Hell-Dunkel-Grenze liegen; damit die Grenze scharf ist, soll es sie
\emph{berühren}.

Das Plugin wählt die Spiegelfläche so, dass genau das gilt: An jedem Punkt einer Facette berührt das
Wendelbild den ihr zugewiesenen Ast der Grenze. Mathematisch ist das eine partielle
Differentialgleichung erster Ordnung für die Fläche. Ihre \textbf{charakteristischen Streifen} --
Kurven auf dem Spiegel mit den Normalen längs der Kurve -- lassen sich als gewöhnliche
Differentialgleichungen integrieren. Längs eines solchen Streifens darf der Spiegel um die Tangente
kippen, ohne dass das Bild die Grenze verlässt; es gleitet nur an ihr entlang.

\paragraph{Eine Facette} ist die Schar der charakteristischen Streifen durch eine Startkurve quer
über den Streifen. Auf der Startkurve ist die Normale ganz vorgegeben: das Bild berührt die Grenze,
und die Bildmitte liegt am \emph{Zielpunkt} (\emph{aim}). Die Mitte der Startkurve liegt auf dem
Grundparaboloid. Auf der ganzen Facette berühren alle Wendelbilder die Grenze -- bis auf die
Genauigkeit der B-Spline-Fläche, typisch unter \ang{0.05} (das Protokoll nennt den Wert). Der
Zielpunkt gilt genau auf der Startkurve und wandert mit dem Abstand davon.

\paragraph{Grundparaboloid.} Alle Facetten beginnen auf einem Paraboloid mit dem Brennpunkt in der
Wendelmitte und der Brennweite \gui{Focal length} (Gruppe \gui{Reflector}, Standard 25 mm).

\subsection{Streifen, Facetten und Reihen}\label{sec:al-streifen}

Der Spiegel wird in \textbf{Zonen} geteilt (Sektoren um die Lampenachse, \autoref{sec:al-zonen}),
jede Zone in \textbf{Streifen} der Breite \gui{Strip width} (Standard 10 mm). Jeder Streifen ist eine
Facette.

\begin{description}
  \item[\gui{Facet length}] Standard 0: ganze Streifen. Ein langer Streifen mit großer Streuung
    verliert zum Rand hin an Streuung, weil der Zielpunkt nur auf der Startkurve exakt gilt. Mit einer
    Facettenlänge größer null wird jeder Streifen in \textbf{Reihen} dieser Länge zerschnitten (quer
    zum Streifen; die Reihe 0 liegt um die Höhe der Lampenachse). Jede Reihe ist eine eigene Facette
    mit eigenen Zielpunkten auf ihrer eigenen Startkurve; jede Reihe beginnt, wo die vorige endet, so
    dass nur eine niedrige Stufe und ein Knick bleiben. Der Preis: mehr Verbindungsflächen.
  \item[\gui{Surface knots every}] Knotenabstand der B-Spline-Flächen der Facetten (Standard
    1,5 mm). Kleiner: näher an den berechneten Streifen, mehr Kontrollpunkte.
\end{description}

Die \textbf{Streifengrenzen} folgen den Charakteristiken einer Facette, die auf die Mitte der
Zielpunkte ihrer Zone zielt: Schnitte des Grundparaboloids mit parallelen Ebenen -- in der
Vorderansicht Geraden für eine geradeaus zielende Zone, sonst Kreise.

Die Facetten heißen im Modell \enquote{Facet <Zone>.<Streifen>} (mit Reihen
\enquote{Facet <Zone>.<Streifen>.<Reihe>}), die Zonen von 1 an gezählt.

\subsection{Zonen}\label{sec:al-zonen}
\index{Zone (Reflektor)}

Die Tabelle \gui{Zones} (Gruppe \gui{Reflector}) teilt den Spiegel in Sektoren um die Lampenachse.
Jede Zeile ist eine Zone mit den Spalten:

\begin{description}
  \item[\gui{from azimuth}] Beginn der Zone (Azimut in Grad, \autoref{sec:al-koordinaten}). Jede Zone
    reicht bis zum Beginn der nächsten (nach Azimut sortiert), die letzte bis zur ersten.
  \item[\gui{cut-off branch}] der Ast, an den die Zone ihre Wendelbilder legt: \code{horizontal}
    oder \code{rising}.
  \item[\gui{aim, right border}] und \gui{aim, left border} der Zielpunkt H (positiv nach rechts) des
    Wendelbilds der Wendelmitte am rechten bzw. linken Rand jedes Streifens. Dazwischen läuft er
    gleichmäßig -- das ist die \textbf{Streuung} des Streifens.
  \item[\gui{below cut-off}] hält die Wendelbilder der Zone so viel unter der Grenze (Standard 0).
\end{description}

\paragraph{Die Standardzonen.} \gui{Default zones} füllt die Tabelle für die gewählte Verkehrsseite.
Für Rechtsverkehr:

\begin{center}
\begin{tabularx}{\linewidth}{rlrrr>{\raggedright\arraybackslash}X}
\toprule
Beginn (\gui{from azimuth}) & Ast & Ziel rechts & Ziel links & \gui{below} &
  Rolle \\
\midrule
$-20$ & horizontal & 6 & $-1$ & 0 & linke Flanke: waagerechte Grenze und Hotspot \\
20 & horizontal & 25 & $-25$ & 0 & oben: Breitenstreuung \\
145 & rising & 7 & 1 & 0 & rechte Flanke: trägt den ansteigenden Ast \\
185 & horizontal & 25 & $-25$ & 0 & unten: Breitenstreuung \\
\bottomrule
\end{tabularx}
\end{center}

Die Flanke, deren Wendelbilder längs des ansteigenden Asts liegen, trägt ihn; die andere Flanke
bildet die waagerechte Grenze und den Hotspot; oben und unten verbreitern das Licht. Für Linksverkehr
werden die Zonen gespiegelt (der Anhang nennt die Werte, \autoref{sec:al-ref-zonen}).

\begin{hinweis}
Zwei Zonen dürfen nicht am selben Azimut beginnen, und die Tabelle darf nicht leer sein -- sonst
meldet das Plugin \enquote{Two zones start at the same azimuth.} bzw. \enquote{The zone table is empty:
add a zone (or press Default zones).}
\end{hinweis}

\paragraph{Wie weit eine Zone zielen kann.} Ein steiler Teil des Spiegels kann sein Licht nicht über
seine eigene Wand werfen. Verlangt der Zielpunkt einer Zone das, wird er zurückgenommen; das
Protokoll nennt die betroffenen Facetten (\enquote{aim held back in … facets}). Zielpunkte, die zu
stark zusammenlaufen, lassen die Streifen einer Facette sich kreuzen -- das Plugin bricht dann mit
einer Meldung ab (\autoref{sec:al-fehler}). Die Startkurve liegt am tiefsten Schnitt der Facette,
damit Streifen mit auseinanderlaufenden Zielen (rechter Rand nach rechts) zum Rand hin auffächern.

\subsection{Wendel quer zur Achse}

Bei einer Querwendel (H3, Signallampen) liegt das Bild an einer waagerechten Grenze, wenn die
Normale des Spiegels keinen Anteil längs der Wendel hat. Eine solche Linie läuft quer über die
Streifen auf Höhe der Lampenachse (die Streifen wechseln dort das berührende Wendelende), eine zweite
längs der Streifen, wo eine Facette seitlich geradeaus blickt. Deshalb liegt bei einer Querwendel eine
Streifengrenze an der Lampenachse (nicht eine Streifenmitte), und jede Facette bleibt auf einer Seite
dieser Linie; wo ihr Zielpunkt sie hinüberbringen würde, wird er verschoben (das Protokoll nennt
diese Facetten).

\begin{tipp}
Für einen Nebelscheinwerfer mit H3: \gui{Rise} = 0 und eine Zone (oder eine obere und eine untere).
Die asymmetrische Grenze mit ansteigender Zone funktioniert ebenfalls.
\end{tipp}

% ---------------------------------------------------------------------------------------------
\section{Verbindungsflächen}\label{sec:al-verbindung}
\index{Verbindungsfläche}\index{Entformungsschräge}

Benachbarte Facetten treffen mit einer Stufe aufeinander, die sich längs ihrer Grenze ändert. Damit
der Spiegel eine Fläche wird, die sich spritzgießen lässt (Entformung entlang X), verbindet das Plugin
die Facetten durch \textbf{Wände mit Entformungsschräge}, an Fuß und Kopf \textbf{verrundet}.

\begin{description}
  \item[\gui{Connecting faces}] an (Standard): eine zusammenhängende Spiegelfläche; aus: nur die
    nackten Facetten mit Stufen.
  \item[\gui{Draft angle}] Winkel zwischen den Wänden und der Entformungsrichtung X (Standard
    \ang{5}, 0,5 bis \ang{45}).
  \item[\gui{Fillet radius}] Radius der Verrundung an Fuß und Kopf jeder Wand (Standard 0,3 mm).
\end{description}

Quer zu einer Grenze legt das Plugin die Wand durch die Mitte der Stufe mit der gewählten Neigung und
rollt dann eine Kugel mit dem Verrundungsradius über beide Seiten. Wo die Stufe verschwindet,
verschwindet auch die Wand; nur der Knick wird verrundet. Die Profile längs einer Grenze ergeben ein
\textbf{Band}, eine bikubische B-Spline-Fläche, die zu beiden Seiten in die Facetten ausläuft; die
Facetten werden darauf zurückgeschnitten. Es gibt drei Ebenen: Bänder zwischen den Reihen eines
Streifens, zwischen den Streifen einer Zone und entlang der Strahlen zwischen zwei Zonen. Alle Bänder
zwischen Zonen haben dieselbe Breite.

\paragraph{Was Sie erwarten können.}
\begin{itemize}
  \item Die Wände halten die Entformungsschräge ein -- etwas weniger nur, wo zwei Wände an einer
    Kreuzung zusammentreffen. Das Protokoll nennt den kleinsten Winkel.
  \item Bänder und Facetten schließen auf etwa 0,01 mm aneinander an.
  \item Mit den Standardeinstellungen nehmen die Bänder 10 bis 15\,\% der Spiegelfläche ein. Ihr
    Licht geht woandershin; in der Lichtverteilung erscheint es als Streulicht über der Grenze.
\end{itemize}

Niedrigere Stufen bedeuten schmalere Wände: schmalere Streifen, weniger Streuung je Streifen oder
benachbarte Zonen mit ähnlichen Zielen.

\begin{warnung}
Die Bänder zwischen den Zonen treffen sich an der Lampenachse. Das Lampenloch muss groß genug für
sie sein; sonst meldet das Plugin den nötigen Mindestradius (\enquote{the lamp hole (radius … mm) is
too small for the connecting faces between the zones, which meet at the lamp axis: at least … mm}).
Ebenso müssen die Streifen breiter sein als ihre Verbindungsflächen.
\end{warnung}

\subsection{Den Reflektor erzeugen}\label{sec:al-generate}

\gui{Generate reflector} erzeugt
\begin{itemize}
  \item jede Facette als getrimmte NURBS-Fläche (bikubischer B-Spline, mit ihrer Zelle getrimmt),
  \item jede Verbindungsfläche (\enquote{Connection …}),
  \item die Wendel als Zylinder (\enquote{Filament <Lampe>}),
\end{itemize}
in der Gruppe \enquote{Reflector <Lampe>} des Modellbaums. Ein neuer Lauf ersetzt den vorigen Reflektor
(alle Flächen mit dem Attribut \code{al.part}). Die Flächen tragen Attribute:

\begin{center}
\begin{tabular}{ll}
\toprule
Attribut & Werte \\
\midrule
\code{al.part} & \code{facet}, \code{connection}, \code{filament} \\
\code{al.zone}, \code{al.strip}, \code{al.row} & Zone (ab 1), Streifen, Reihe \\
\code{al.branch} & \code{horizontal} oder \code{rising} \\
\bottomrule
\end{tabular}
\end{center}

Das Protokoll nennt Zahl der Facetten und Zonen, Spiegelfläche (Vorderansicht, cm$^2$), Tiefe in X,
wie nah die Wendelbilder an der Grenze liegen, die höchste Stufe, den Verrundungsradius, den
kleinsten Entformungswinkel und den Flächenanteil der Verbindungsflächen. Die Statusleiste etwa:
\enquote{<Anzahl> facets, <Anzahl> connecting faces, images within <Winkel>\si{\degree} of the cut-off}.

Die Flächen sind einzelne getrimmte Flächen, nicht zu einer Schale vernäht.

\subsection{Streifen anzeigen}

\gui{Show strips} zeigt im Ergebnissatz \enquote{Strips <Lampe>} die \enquote{Start curves} und
einige \enquote{Characteristic strips} jeder Facette als Linien in den Ansichten. So sehen Sie, wie die
Facetten aufgebaut sind und wo Streifen auffächern oder sich drängen. Die Aktion gilt nur für den
facettierten Reflektor.

% ---------------------------------------------------------------------------------------------
\section{Paraboloid und Streuscheibe}\label{sec:al-streuscheibe}
\index{Streuscheibe}\index{Paraboloid}

Mit \gui{System} = \gui{patterned lens: paraboloid and lens} besteht der Scheinwerfer aus einem
Paraboloidspiegel und einer Streuscheibe davor. Es gelten die Gruppen \gui{Frame}, \gui{Lamp} und
\gui{Lens}; aus \gui{Reflector} nur \gui{Focal length}.

\subsection{Der Spiegel}

Der Spiegel ist das Paraboloid mit \gui{Focal length} innerhalb des Rahmens, außerhalb des
Lampenlochs. \gui{Mirror focus behind the filament centre} (Gruppe \gui{Lens}, Standard 0) verschiebt
seinen Brennpunkt hinter die Wendelmitte. Bei der H4 sitzt die Abblendwendel vor dem Brennpunkt
(im H4-Scheinwerfer liegt er an der Fernlichtwendel, 6,7 mm dahinter). Dann sendet die obere
Spiegelhälfte -- die allein Licht am Schirm vorbei bekommt -- ihr Licht unter den Horizont.

\subsection{Form der Scheibe}

\begin{description}
  \item[\gui{Lens in front of the filament}] Abstand der Grundebene der Innenseite vor der Wendel,
    auf der Lampenachse (Standard 75 mm).
  \item[\gui{Thickness}] Dicke der Scheibe, längs X gemessen (Standard 3 mm).
  \item[\gui{Rake}] die Scheibe neigt sich zurück: ihr oberer Rand zum Spiegel hin (Standard
    \ang{0}).
  \item[\gui{Sweep}] die Scheibe läuft nach links zurück (negativ: nach rechts).
  \item[\gui{Dome radius}] nach außen gewölbt als Kugel dieses Radius; 0: eben.
  \item[\gui{Refractive index}] Brechzahl (Standard 1,49; PMMA 1,49, Polycarbonat 1,59, Glas etwa
    1,52).
  \item[\gui{Transmittance}] Transmissionsgrad der Scheibe (Standard 0,92) -- bei einer roten oder
    gelben Scheibe einschließlich der Farbe.
\end{description}

Die Außenseite ist glatt und hat die gewählte Form; die Innenseite hat dieselbe Form, \gui{Thickness}
dahinter, und trägt die Elemente.

\subsection{Zellen und Elemente}\label{sec:al-elemente}
\index{Prisma}\index{Zylinderlinse}\index{Kissen}

Die Innenseite ist ein Raster aus Zellen der Größe \gui{Element width} $\times$ \gui{Element height}
(Standard je 10 mm); eine Zelle liegt mittig auf der Lampenachse. Was jede Zelle tut, legt die
Tabelle \gui{Elements} fest. Jede Zeile beschreibt einen Bereich (y nach links, z nach oben, um die
Lampenachse, in mm) und was die Zellen tun, deren Mitte darin liegt; die erste passende Zeile gilt.

\begin{description}
  \item[\gui{y from}, \gui{y to}, \gui{z from}, \gui{z to}] der Bereich (Standard $\pm 500$ mm, also
    alles).
  \item[\gui{deflect H}, \gui{deflect V}] um wie viel die Zelle das Licht in ihrer Mitte ablenkt --
    ein \textbf{Prisma}.
  \item[\gui{spread $\pm$ H}, \gui{spread $\pm$ V}] um wie viel mehr an ihren Rändern, linear über
    die Zelle (der linke Teil nach links, der obere nach oben; negativ: umgekehrt, die Strahlen
    kreuzen sich) -- eine \textbf{Zylinderlinse} mit einer Streuung, ein \textbf{Kissen} mit beiden.
  \item[\gui{spread turned by}] beide Streuungen auf dem Schirm um diesen Winkel gedreht (gegen den
    Uhrzeigersinn, wie der Fahrer sieht). Eine Zylinderlinse (\gui{spread $\pm$ V} = 0), um
    \ang{15} gedreht, streut das Licht längs einer Linie, die um \ang{15} ansteigt -- so wie die
    Facetten eines Klarglasreflektors ihre Wendelbilder längs der Grenze legen. Die Streuung wird an
    der Zellenecke erreicht, die in ihrer Richtung am weitesten außen liegt.
  \item[\gui{sector}] der Sektor, für dessen Teil der Zelle die Zeile gilt (ihre Nummer nach Azimut,
    ab 0); $-1$: jeder Teil.
\end{description}

Alle Werte 0: ebenes Glas. Die Standardzeile ist ein Kissen über die ganze Scheibe:
\gui{deflect V} = $-1$, \gui{spread $\pm$ H} = 15, \gui{spread $\pm$ V} = 4.

\paragraph{Wie ein Element entsteht.} Für einen Punkt der Zelle soll der Strahl, der von der Wendelmitte
über den Spiegel kommt, die Scheibe um die gewählten Winkel gedreht verlassen. Aus der gewünschten
Austrittsrichtung folgt mit der Brechung an der Außenseite die Richtung im Glas, daraus mit dem
Snelliusschen Gesetz in Vektorform die Normale der Innenseite. Diese Neigungen werden über die Zelle
integriert; die Fläche der Zelle ist das B\'ezier-Flächenstück vom Grad 5, das dem am nächsten kommt
(das Protokoll nennt den größten Neigungsfehler). Eine Scheibe ohne Elemente lenkt nichts ab.

\paragraph{Wände zwischen den Elementen.} Benachbarte Elemente treffen mit einer Stufe aufeinander.
Sie wird durch eine Wand längs X geschlossen, als Regelfläche zwischen den Rändern der beiden
Flächenstücke. Ein Strahl, der eine Wand trifft, wird dort einmal gespiegelt und geht als Streulicht
weiter: aus der Luft vor den Elementen mit dem Reflexionsgrad des Glases unter diesem streifenden
Winkel (was in die Wand eindringt, geht verloren), im Glas dahinter durch Totalreflexion.
\gui{Light distribution} meldet, wie viel Licht die Wände trifft und wie viel davon weitergeht.

\paragraph{Wände mit Schräge.} Mit \gui{Draft angle of the walls} größer null neigen sich die Wände
um diesen Winkel gegen X und werden an Fuß und Kopf mit \gui{Fillet radius of the walls} verrundet --
wie es Spritzgießen oder Pressen verlangen. Sie entstehen wie die Verbindungsflächen des
Klarglasreflektors als Bänder, die in die Elemente auslaufen (zuerst zwischen den Elementen einer
Reihe, dann zwischen den Reihen). Die Wände sind dann Teil der Innenfläche: ein Strahl wird an ihnen
gebrochen wie überall sonst.

\begin{warnung}
Wände mit Schräge und Sektoren lassen sich noch nicht kombinieren. Setzen Sie die Schräge auf 0 oder
lassen Sie die Sektortabelle leer.
\end{warnung}

\subsection{Sektoren}\label{sec:al-sektoren}
\index{Sektor (Streuscheibe)}

Die Tabelle \gui{Sectors} teilt die Scheibe in radiale Abschnitte um die Lampenachse -- wie die Zonen
eines Klarglasreflektors. Jede Zeile beginnt bei \gui{from azimuth} (0 links, 90 oben) und reicht bis
zum Beginn der nächsten; \gui{cut-off branch} nennt den Ast, an dem das Licht des Sektors endet.

Mit zwei oder mehr Sektoren werden Zellen, die eine radiale Linie schneidet, entlang dieser Linie
geteilt: jeder Teil ist ein eigenes Element (sein Flächenstück über die ganze Zelle entworfen, seine
Ablenkung um die Mitte des Teils, seine Streuung über die Ausdehnung des Teils). Die Stufe zwischen
zwei Teilen wird eine Wand entlang der radialen Linie. In der Elementtabelle wählt die Spalte
\gui{sector} den Teil.

\paragraph{Warum Sektoren?} Das Licht der H4, das den Schirm über seine schräge Kante verlässt,
erreicht einen Keil aus der Mitte der Scheibe, dessen Ränder radiale Linien sind -- eine klassische
H4-Scheibe hat dort ihre schrägen Riefen. Zellen eines quadratischen Rasters mischen es an beiden
Rändern mit dem übrigen Licht: etwa 40\,\% dieses Lichts liegen in Zellen, die halb drin, halb draußen
sind. Leer (Standard) gilt ein einziger Sektor; \gui{Optimise lens} mit \gui{Optimise: along the
edges} füllt die Tabelle dann selbst (\autoref{sec:al-optimieren}).

\subsection{Den Scheinwerfer erzeugen}

\gui{Generate reflector} erzeugt bei dieser Bauart in der Gruppe \enquote{Headlamp <Lampe>}:
\begin{itemize}
  \item den Spiegel, exakt, mit Lampenloch (\enquote{Mirror (paraboloid)}, \code{al.part} =
    \code{mirror}),
  \item die Außenseite der Scheibe (\enquote{Lens outer side}, \code{lens}),
  \item jedes Element (\enquote{Lens element i.j}, \code{element}) und jede Wand (\code{wall}),
  \item den Wendelzylinder (\code{filament}).
\end{itemize}

Das Protokoll nennt Brennweite, Lage, Form, Dicke und Brechzahl der Scheibe, die Zahl der Elemente,
die Tiefe ihrer Flächen und den Neigungsfehler; die Statusleiste \enquote{paraboloid and lens with …
elements}.

\begin{hinweis}
Mögliche Abbrüche beim Aufbau der Scheibe: Die Scheibe liegt am Rand im Spiegel (weiter nach vorn
oder weniger neigen), ihr Kugelradius ist kleiner als ihre Breite, ein Element kann das Licht nicht so
weit lenken, die Elemente sind tiefer als die Scheibe dick ist (dickere Scheibe, kleinere Zellen,
weniger Ablenkung), oder sie sind für Wände mit dieser Schräge zu klein. Die Meldungen stehen in
\autoref{sec:al-fehler}.
\end{hinweis}

% ---------------------------------------------------------------------------------------------
\section{Wendelbilder}\label{sec:al-wendelbilder}

\gui{Filament images} zeigt im Fenster \gui{Image} (Ergebnissatz \enquote{Filament images <Lampe>}) die
Umrisse der Wendelbilder von Punkten über den ganzen Spiegel, zusammen mit der Hell-Dunkel-Grenze als
heller Linie. Der Bereich ist H $= -45$ bis $45^\circ$, V $= -15$ bis $5^\circ$, 20 Bildpunkte je
Grad; jeder Umriss zählt 1, überlappende Umrisse addieren sich.

\gui{Filament images (number)} (Gruppe \gui{Analysis}, Standard 150) legt fest, wie viele Spiegelpunkte
ungefähr genommen werden. Beim Klarglasreflektor sind es Punkte der Facetten, bei der Streuscheibe
Punkte des Spiegels, deren Licht durch die Scheibe verfolgt wird.

Damit prüfen Sie auf einen Blick:
\begin{itemize}
  \item ob alle Bilder unter der Grenze liegen und sie berühren,
  \item wie die Zonen das Licht verteilen (Hotspot, Breite, ansteigender Ast),
  \item ob eine Zone zu weit seitlich zielt (Bilder, die sich am Rand stauen).
\end{itemize}

% ---------------------------------------------------------------------------------------------
\section{Lichtverteilung und Prüfpunkte}\label{sec:al-lichtverteilung}
\index{Lichtverteilung}\index{UN/ECE R112}\index{Prüfpunkt}

\gui{Light distribution} berechnet die Lichtstärke des Scheinwerfers auf dem Messschirm.

\subsection{Das Modell}

Die Wendel ist ein Zylinder, dessen Mantel in alle Richtungen gleich hell leuchtet (Lambertsch); jeder
Spiegelpunkt -- auf den Facetten und den Verbindungsflächen -- reflektiert einmal, mit dem
Reflexionsgrad aus der Gruppe \gui{Analysis}. Strahlen gehen von jedem Abtastpunkt des Wendelmantels zu
jedem Spiegelpunkt eines Gitters in der Vorderansicht. Bei der Streuscheibe folgt die Brechung an
beiden Seiten der Scheibe mit \gui{Transmittance}; ein Strahl, der die Scheibe verfehlt oder an ihrer
Außenseite total reflektiert wird, ist verloren, einer, der eine Wand trifft, wird dort gespiegelt.

Die Einstellungen der Gruppe \gui{Analysis}:

\begin{description}
  \item[\gui{Reflectance}] Reflexionsgrad des Spiegels (Standard 0,85).
  \item[\gui{Reflectance of the connecting faces}] Reflexionsgrad der Verbindungsflächen (Standard
    0,85); 0: ihr Licht zählt nicht (abgedeckte oder matte Wände).
  \item[\gui{Test points}] welche Prüfpunkte die Tabelle zeigt: \gui{low beam (UN R112, class B)}
    (Standard) oder eine Signalleuchte nach UN R148 (\autoref{sec:al-signal}).
  \item[\gui{Shadows of the steps}] (Standard an) beim facettierten Reflektor mit Verbindungsflächen:
    ein Strahl, den eine Stufe des Spiegels auf dem Weg von der Wendel aufhält, entfällt; einer, den
    sie nach der Reflexion aufhält, wird dort ein zweites Mal reflektiert (mit dem Reflexionsgrad der
    Verbindungsflächen) und zählt als Streulicht, wenn er den Spiegel danach frei verlässt. Kostet
    Rechenzeit.
  \item[\gui{Direct light of the lamp}] (Standard an) das Licht, das die Wendel nach vorn am Spiegel
    vorbei verlässt -- über den Schirm einer H4 und um ihre Kappe herum, bei der Streuscheibe durch
    die Scheibe.
  \item[\gui{Mirror points every}] Gitterabstand der Spiegelpunkte (Standard 1 mm, 0,2 bis 10 mm).
    Kleiner: genauer und langsamer.
\end{description}

\subsection{Die Ergebnisse}

Im Ergebnissatz \enquote{Low beam <Lampe>} erscheinen

\begin{description}
  \item[Bild \enquote{Luminous intensity}] die Lichtstärke in cd über H und V (Fenster \gui{Image}).
    Der Bereich ist H $= -40$ bis $40^\circ$, V $= -15$ bis $10^\circ$ in Zellen von \ang{0.25}
    $\times$ \ang{0.1}; bei Signalleuchten H $= \pm 40^\circ$, V $= \pm 20^\circ$ in Zellen von
    \ang{0.25} $\times$ \ang{0.25}.
  \item[Tabelle der Prüfpunkte] (Fenster \gui{Table}, siehe unten).
  \item[Kurve \enquote{Vertical cut at H = …}] ein senkrechter Schnitt durch die Grenze bei
    H $= -2{,}5^\circ$ (Rechtsverkehr; Linksverkehr $+2{,}5^\circ$), gemittelt über $\pm 0{,}5^\circ$.
  \item[Kurve \enquote{Horizontal cut at V = -0.86}] ein waagerechter Schnitt in der Höhe der
    50-m-Linie, gemittelt über $\pm 0{,}25^\circ$. Bei Signalleuchten gehen beide Schnitte durch die
    Achse.
\end{description}

Im Protokoll stehen
\begin{itemize}
  \item die Lichtströme: der Lampe, auf den Spiegel, im Bündel und davon aus den Verbindungsflächen
    (Streulicht);
  \item die höchste Lichtstärke;
  \item die Lage der Grenze und ihre \textbf{Schärfe} $G$ (siehe unten);
  \item die höchste Lichtstärke mehr als \ang{0.2} über der Grenze (Blendung);
  \item mit \gui{Shadows of the steps} der Anteil, den Stufen aufhalten;
  \item mit \gui{Direct light of the lamp} dessen Lichtstrom und Höchstwert (getrennt vom Bündel);
  \item bei der Streuscheibe der Lichtstrom, der die Wände trifft, und wie viel davon als Streulicht
    weitergeht.
\end{itemize}

Die Statusleiste fasst zusammen, etwa \enquote{… lm in the beam (… \%), maximum … cd, cut-off G =
…}.

\subsection{Die Prüfpunkte des Abblendlichts}

Die Tabelle \enquote{R112 test points (class B, passing beam)} hat die Spalten \code{point},
\code{H / }\si{\degree}, \code{V / }\si{\degree}, \code{I / cd}, \code{E at 25 m / lx},
\code{required} und das Urteil (\code{ok} oder \code{not met}). Sie enthält die Punkte nach UN/ECE
R112, Absatz 6.2.4 (Scheinwerfer der Klasse B, Abblendlicht, Rechtsverkehr; für Linksverkehr an
V-V gespiegelt):

\begin{center}
\begin{tabular}{lrrl}
\toprule
Punkt & H / \si{\degree} & V / \si{\degree} & Anforderung \\
\midrule
B50L & $-3{,}43$ & 0,57 & $\leq 350$ cd \\
BR & 2,5 & 1,0 & $\leq 1750$ cd \\
75R & 1,15 & $-0{,}57$ & $\geq 10\,100$ cd \\
75L & $-3{,}43$ & $-0{,}57$ & $\leq 10\,600$ cd \\
50L & $-3{,}43$ & $-0{,}86$ & $\leq 13\,200$ cd \\
50R & 1,72 & $-0{,}86$ & $\geq 10\,100$ cd \\
50V & 0 & $-0{,}86$ & $\geq 5100$ cd \\
25L & $-9$ & $-1{,}72$ & $\geq 1700$ cd \\
25R & 9 & $-1{,}72$ & $\geq 1700$ cd \\
Zone III & \multicolumn{2}{l}{Vieleck über der Grenze} & höchstens 625 cd \\
Zone IV & \multicolumn{2}{l}{H $\pm 5{,}15$, V $-1{,}72$ bis $-0{,}86$} & überall $\geq 2500$ cd \\
Zone I & \multicolumn{2}{l}{H $\pm 9$, V $-4$ bis $-1{,}72$} & höchstens $2 \times I(\text{50R})$ \\
\bottomrule
\end{tabular}
\end{center}

Ein Prüfpunkt ist der Mittelwert über das Fenster eines Photometers ($\pm \ang{0.25}$ in H,
$\pm \ang{0.1}$ in V). Die Zonen III und I prüfen den höchsten, Zone IV den kleinsten solchen
Mittelwert.

\begin{warnung}
Die Tabelle ist eine Entwurfshilfe. Maßgeblich ist die Regelung selbst, nicht diese Tabelle.
\end{warnung}

\paragraph{Schärfe der Grenze.} Wie R112 misst das Plugin die Schärfe auf der senkrechten Linie
\ang{2.5} zur Seite des Gegenverkehrs:
$G = \log I(V) - \log I(V + \ang{0.1})$, der größte Wert innerhalb \ang{2} der Nenngrenze, nur wo
$I(V)$ mindestens 2\,\% des Linienmaximums ist. R112 verlangt 0,13 bis 0,40. Das Modell rechnet die
Wendel als gleichmäßig leuchtenden Zylinder und kennt keine Streuung; eine echte Lampe ist weicher,
$G$ fällt im Modell deshalb höher aus. Das Raytracing kommt der Wirklichkeit näher
(\autoref{sec:al-raytracing}).

\subsection{Signalleuchten}\label{sec:al-signal}
\index{Signalleuchte}\index{UN R148}

Mit \gui{Test points} auf einer Signalleuchte zeigt \gui{Light distribution} die Punkte ihrer
Standard-Lichtverteilung nach UN R148 ($\pm \ang{20}$ horizontal, $\pm \ang{10}$ vertikal: der
geforderte Prozentsatz der Mindestlichtstärke auf der Achse) und die höchste Lichtstärke in jeder
Richtung -- für eine einzelne Leuchte:

\begin{center}
\begin{tabular}{lrr}
\toprule
\gui{Test points} & Minimum auf der Achse & Maximum \\
\midrule
\code{front position lamp (A)} & 4 cd & 140 cd \\
\code{rear position lamp (R1)} & 4 cd & 17 cd \\
\code{stop lamp (S1)} & 60 cd & 260 cd \\
\code{direction indicator, front (1)} & 175 cd & 1000 cd \\
\code{direction indicator, rear (2a)} & 50 cd & 500 cd \\
\code{daytime running lamp (RL)} & 400 cd & 1200 cd \\
\bottomrule
\end{tabular}
\end{center}

Die Tabelle heißt dann \enquote{R148 standard light distribution, …} mit den Spalten \code{point},
\code{H / }\si{\degree}, \code{V / }\si{\degree}, \code{I / cd}, \code{required} und Urteil; ein Punkt
ist der Mittelwert über $\pm \ang{0.5}$. Die Tagfahrleuchte hat ihr eigenes Punktraster. Die Farbe der
Scheibe ist Teil von \gui{Transmittance}: eine rote oder gelbe Scheibe lässt nur einen Teil des
Lichtstroms durch. Die Lampen P21W, PY21W, R10W und R5W haben eine Querwendel.

\subsection{Direktlicht der Lampe}

Mit \gui{Direct light of the lamp} geht das Licht, das die Wendel nach vorn am Spiegel vorbei
verlässt, in Bild und Tabelle ein; sein Lichtstrom steht getrennt im Protokoll. Die Kappe einer H4
hält es innerhalb von etwa \ang{40} um die Achse auf; eine Lampe ohne Kappe verteilt es über den
ganzen Schirm -- bei einer H7 bis etwa 100 cd. Das Direktlicht wird auf einem Raster von \ang{1}
genommen und, ob der Spiegel im Weg ist, von der Wendelmitte aus beurteilt.

% ---------------------------------------------------------------------------------------------
\section{Raytracing mit der Lampe, wie sie ist}\label{sec:al-raytracing}
\index{Raytracing (Automotive)}\index{Path Tracer}

\gui{Light distribution (ray tracing)} prüft den Entwurf ohne die Annahmen der Wendelbilder. Die
Wendelbilder bleiben für Entwurf, \gui{Filament images} und \gui{Light distribution} unverändert;
das Raytracing ist die Kontrolle.

\subsection{Was das Raytracing nachbildet}

\begin{description}
  \item[Die Wendel als Helix.] Ein Wolframdraht füllt den Zylinder der Lampe (Länge, Durchmesser):
    die Drahtmitte läuft auf einer Helix mit dem Radius (Durchmesser $-$ Draht)/2, je Vierteldrehung
    eine NURBS-Fläche (ein Rohr um die Helix, auf \qty{0.2}{\micro\metre} genau). Die Oberfläche
    strahlt Lambertsch; Licht, das eine andere Windung trifft, wird absorbiert. Der Lichtstrom ist so
    gewählt, dass das, was aus der Wendel herauskommt, der Lichtstrom der Lampe ist -- bei einer H7
    kommen etwa 56\,\% des Drahtlichts heraus.
  \item[Drahtdurchmesser und Windungen.] Sie stehen nicht in R37. Ohne eigene Werte
    (\gui{Coil: wire diameter} und \gui{Coil: turns} = 0) schätzt das Plugin sie aus der Leistung:
    ein Draht bei \qty{3100}{K} bei 13,2 V, gewickelt mit einer Steigung von 1,5 Drahtdurchmessern.
    Das ergibt bei der H7 \qty{0.161}{mm} und 16 Windungen, bei der H4 \qty{0.187}{mm} und 19.
  \item[Kühlere Enden.] Die Zuleitungen ziehen Wärme von den Wendelenden ab. Aus der
    Wärmebilanz längs des Drahts steigt die Temperatur von \gui{Coil: temperature at its ends}
    (angenommen \qty{2400}{K}) zur Mitte (\qty{3100}{K}): nach 2 mm Draht etwa \qty{2900}{K}, nach
    4 mm etwa \qty{3040}{K}. Jede Vierteldrehung leuchtet mit der Leuchtdichte ihrer Temperatur
    (Plancksches Gesetz, mit $V(\lambda)$ gewichtet); die erste und letzte Windung einer H7 etwa mit
    der Hälfte der Mitte. Mit \qty{3100}{K} leuchtet die Wendel gleichmäßig.
  \item[Der Kolben.] Ein Glasrohr um die Bezugsachse mit einer Halbkugel als Spitze, ein Glaskörper
    mit Fresnelscher Reflexion und Brechung (Quarz $n = 1{,}458$, Glas der Signallampen 1,52), hinten
    offen. Seine schwarze Kuppe -- eine absorbierende Schicht außen auf dem Glas -- reicht so weit
    zurück, wie R37 mindestens verlangt. Durchmesser, Wand und Lage der Spitze sind typische Werte
    (H7 9 mm, H1 8,5 mm, H4 13 mm, Wand 1 mm; P21W 26,5 mm, R5W/R10W 19 mm). Ohne Glas
    (\gui{Bulb glass} aus) ist die Kuppe eine Scheibe quer über den Kolben.
  \item[Das Innere der H4.] Der Schirm (Wanne und Stirnwand) als Spiegel mit
    \gui{H4 shield: reflectance} (Standard 0,5; 0: schwarz) und die zweite Wendel, dunkel, an ihrem
    Platz (4,5 mm lang, 1,7 mm hinter der Abblendwendel, ihre Achse 0,8 mm unter der Ebene H).
  \item[Spiegel und Scheibe.] Jede Facette und Verbindungsfläche (oder das Paraboloid) wie erzeugt,
    mit den Reflexionsgraden des Panels; die Scheibe als ein Glaskörper mit Fresnelscher Reflexion und
    Brechung bei der Brechzahl des Panels -- \gui{Transmittance} wird hier nicht verwendet, es gibt
    keine Absorption im Glas.
\end{description}

\subsection{Lichtwege und Rechenkern}

Lichtwege beginnen auf dem Draht und werden durch jede Reflexion und Brechung verfolgt, bis zu
\gui{Longest path} Segmenten (Standard 20; allein durch den Kolben sind es zwei). Das austretende Licht
wird nach seiner Richtung in die Zellen des Schirms sortiert. Licht, dessen erste Fläche nach der
Lampe (Wendeln, Kolben, Schirm, Kuppe) der Spiegel ist, ist das Bündel; der Rest, der nach vorn
austritt, ist das Direktlicht. Mit Streuscheibe zählt nur, was durch die Scheibe austritt; der Rest
bleibt im Gehäuse.

\begin{description}
  \item[\gui{Light paths}] Zahl der Lichtwege (Standard 5 Millionen, 10\,000 bis 2 Milliarden). Das
    Rauschen einer Schirmzelle fällt mit der Wurzel: bei 5 Millionen etwa 5\,\% dort, wo das Bündel
    am hellsten ist; vierfach so viele Wege halbieren es.
  \item[\gui{Kernel}] \code{automatic} (Standard: C++, wenn installiert, sonst NumPy), \code{C++} oder
    \code{NumPy}. \code{C++} ohne installiertes Paket \code{pathtracer} meldet einen Fehler mit dem
    Installationsbefehl.
\end{description}

\subsection{Die Ergebnisse}

Der Ergebnissatz \enquote{Low beam <Lampe> – ray tracing} steht neben dem der Wendelbilder, mit
demselben Bild, derselben Tabelle und denselben Schnitten. Das Protokoll nennt
\begin{itemize}
  \item den Rechenkern, die Zahl der Lichtwege, Strahlen und Dreiecke und die Zeit,
  \item die Wendel: Draht, Windungen, Steigung, Drahtlänge (geschätzte Werte als solche markiert) und
    wie viel des Drahtlichts aus der Wendel herauskommt,
  \item die Leuchtdichte der kühleren Enden,
  \item den Kolben (Glas, Durchmesser, Wand, Spitze, Beginn der Kuppe), den Schirm der H4,
  \item die Lichtströme: im Bündel, davon über die Verbindungsflächen, direkt, nach hinten, ins
    Gehäuse (mit Scheibe) und absorbiert,
  \item das Rauschen (Standardfehler) dort, wo das Licht am hellsten ist.
\end{itemize}

\subsection{Fernlicht der H4}

Mit \gui{H4: lit coil} = \gui{driving beam} leuchtet statt der Abblendwendel die Fernlichtwendel
(1650 lm) -- das Fernlicht desselben Reflektors. Der Ergebnissatz heißt dann \enquote{Driving beam H4 –
ray tracing}, die Tabelle \enquote{R112 test points (class B, driving beam)} mit:

\begin{center}
\begin{tabular}{lrrl}
\toprule
Punkt & H / \si{\degree} & V / \si{\degree} & Anforderung \\
\midrule
IM (Maximum) & & & 40\,500 bis 215\,000 cd \\
HV & 0 & 0 & $\geq 80\,\%$ von IM \\
H-5L & $-5$ & 0 & $\geq 5100$ cd \\
H-2.5L & $-2{,}5$ & 0 & $\geq 20\,300$ cd \\
H-2.5R & 2,5 & 0 & $\geq 20\,300$ cd \\
H-5R & 5 & 0 & $\geq 5100$ cd \\
\bottomrule
\end{tabular}
\end{center}

Nur die H4 hat eine Fernlichtwendel; bei anderen Lampen meldet das Plugin \enquote{Only the H4 has a
driving-beam coil (H4: lit coil).}

\subsection{Die Lampe ins Modell legen: Generate lamp}

\gui{Generate lamp} legt die Lampe so ins Modell, wie das Raytracing sie sieht -- zum Ansehen, was
gerechnet wird. In der Gruppe \enquote{Lamp <Lampe>} am \gui{Filament centre} entstehen Untergruppen:

\begin{center}
\begin{tabularx}{\linewidth}{ll>{\raggedright\arraybackslash}X}
\toprule
Gruppe & \code{al.lamp} & Flächen \\
\midrule
\enquote{Coil (lit)} & \code{coil} & der Draht in Vierteldrehungen, benannt \enquote{Coil 3.2}
  (dritte Windung, zweites Viertel) \\
\enquote{Other coil (dark)} & \code{other coil} & die zweite Wendel der H4 \\
\enquote{Shield} & \code{shield} & Wanne und Stirnwand der H4 \\
\enquote{Bulb glass} & \code{bulb} & der Kolben \\
\enquote{Black top} & \code{cap} & die schwarze Kuppe bzw. Kappe \\
\bottomrule
\end{tabularx}
\end{center}

Jede Fläche trägt \code{al.lamp} und \code{al.material}; die Drahtstücke zusätzlich
\code{al.luminance}, ihre Leuchtdichte relativ zur Wendelmitte. Es gelten die Einstellungen der
Gruppe \gui{Ray tracing} (Draht, Windungen, Kolben, leuchtende Wendel). Ein neuer Lauf ersetzt die
vorige Lampe; \gui{Generate reflector} lässt sie stehen.

\begin{tipp}
Glas wird wie jede andere Fläche gezeichnet. Um die Wendel zu sehen, blenden Sie im Modellbaum die
Gruppe \enquote{Bulb glass} aus -- und den Wendelzylinder von \gui{Generate reflector}, den der Draht
von innen berührt. Das Protokoll erinnert daran.
\end{tipp}

Das Raytracing liest diese Flächen nicht; es baut die Lampe jedes Mal neu aus dem Panel.

\subsection{Wendelbilder und Raytracing im Vergleich}

Laut den Vergleichsrechnungen der Entwickler (H7, Testreflektor, 12 Millionen Wege) stimmen
Lichtstrom im Bündel (508 gegen 500 lm), Direktlicht (736 gegen 740 lm) und Blendung über der Grenze
(288 gegen 293 cd) überein; das Maximum ist im Raytracing etwa 15\,\% niedriger, weil der Draht quer
zur Wendel weniger hell leuchtet als der Zylinder und mehr längs ihrer Achse.

Was die echte Lampe ändert (Standardwerte, H7, 5 Millionen Wege, gegen die gleichmäßige Wendel ohne
Glas): die kühleren Enden wenig; der Kolben am meisten -- seine Reflexionen senden Licht über die
Grenze (Zone III 520 statt 200 cd, B50L 180 statt 15 cd, $G = 0{,}36$): Geisterbilder der Wendel, die
R37 die Lampenhersteller über der Wendel zu halten bittet. Bei der H4 ähnlich ($G$ 0,25 statt 0,97,
B50L 93 statt 20 cd).

% ---------------------------------------------------------------------------------------------
\section{Die Streuscheibe optimieren}\label{sec:al-optimieren}
\index{Optimierung (Streuscheibe)}

Ein Abblendlicht mit Streuscheibe bekommt seine Grenze von der Lampe -- dem Schirm der H4, mit dem
Brennpunkt des Paraboloids hinter der Wendel. Die Scheibe muss das Licht unter der Grenze fein
verteilen. \gui{Optimise lens} gibt dafür jeder Zelle ihr eigenes Prisma und ihre eigene
Zylinderlinse -- Stärke und Richtung des Zylinders --, so dass eine Zelle ihr Licht längs der
Grenze legen kann: waagerecht unter den waagerechten Ast, unter \ang{15} unter den Anstieg, steiler
ins Vorfeld.

\paragraph{Voraussetzungen.}
\begin{itemize}
  \item \gui{System} = \gui{patterned lens: paraboloid and lens} (sonst: \enquote{Optimise lens
    belongs to the patterned lens (System: patterned lens).}),
  \item \gui{Test points} = \gui{low beam (UN R112, class B)} (sonst: \enquote{Optimise lens aims at
    the low beam's test points (Analysis: test points).}),
  \item eine Lampe mit Schirm (H4). Ohne Schirm macht das Paraboloid keine Grenze; das Plugin warnt im
    Protokoll, dass die Scheibe allein die Grenzwerte über der Grenze nicht einhalten kann,
  \item für \gui{Optimise: the lamp's blur} = \gui{measured by ray tracing} ein verfügbarer Rechenkern
    (\autoref{sec:al-raytracing}).
\end{itemize}

Kleine Zellen geben der Optimierung mehr Freiheit -- am besten einige Millimeter (\gui{Element width},
\gui{Element height}).

\subsection{Wie die Optimierung arbeitet}

\begin{description}
  \item[Das Licht der Zellen.] Strahlen von der Wendeloberfläche über die Spiegelpunkte erreichen die
    Zellen einer ebenen Scheibe; das Licht jeder Zelle, auf ein Raster von $\ang{0.5} \times
    \ang{0.1}$ sortiert, ist ihr \emph{Muster}. Ein Element verschiebt es um seine Ablenkung und
    verschmiert es längs der Zylinderrichtung über $\pm$ seine Streuung. Das Bündel einer beliebigen
    Wahl ist die Summe der verschobenen und verschmierten Muster plus dem Direktlicht -- ohne die
    Wendelbilder neu zu nehmen.
  \item[Die Auswahl.] Ablenkung $-16$ bis $16^\circ$ in Schritten von \ang{1} quer, $-3$ bis
    $0{,}6^\circ$ in Schritten von \ang{0.2} auf und ab; ein Zylinder mit 0 bis \ang{12} Streuung,
    gedreht um $-60$ bis $90^\circ$ (0, $\pm 5$, 10, 15, 20, 30, 45, 60, 90); $|$Ablenkung$|$ plus
    der Queranteil der Streuung höchstens \ang{22}, ihr Höhenanteil höchstens \ang{4}; nur Elemente,
    die in die Scheibendicke passen.
  \item[Das Ziel.] Die Prüfpunkte von R112 (ihre Photometerfenster) und die Zonen III, IV und I mit
    der Reserve \gui{Optimise: margin} (Standard 10\,\%), als quadrierte Unter- und
    Überschreitungen; danach wird die kleinste Reserve so groß wie möglich. Zwei Forderungen über
    R112 hinaus halten die Grenze, wo sie hingehört: ab \ang{0.3} über der Linie ($|H| \leq 10^\circ$)
    höchstens die 625 cd von Zone III, und bei \ang{2.5} zur Seite des Gegenverkehrs ein Abfall um den
    Faktor 10 von knapp unter zu knapp über der Linie. Wo nichts von einer Zelle abhängt, werden
    kleine Elemente bevorzugt.
  \item[Ein echtes Bündel.] Was R112 nicht verlangt, ein Scheinwerfer aber hat -- mit einem Fünftel
    des Gewichts von R112, damit R112 vorgeht: am hellsten an der Grenze, nicht in einem Fleck darunter
    (\gui{Optimise: under the cut-off}, Standard 6000 cd: mindestens das im Mittel 0,15 bis
    $0{,}65^\circ$ unter beiden Ästen, $\pm \ang{10}$ um den Knick), und der Knick sichtbar mit
    hellem, scharfem Anstieg (\gui{Optimise: under the rise}, Standard 3000 cd: mindestens das 0,15
    bis $0{,}35^\circ$ unter dem ansteigenden Ast, 1 bis $5^\circ$ vom Knick, darüber höchstens ein
    Zehntel). 0 lässt die jeweilige Forderung weg.
  \item[Gleichmäßigkeit.] \gui{Optimise: evenness} (Standard 1, ein Gewicht; 0: nur R112): Unter der
    Grenze (von \ang{0.4} unter der Linie bis $-5^\circ$, $|H| \leq 20^\circ$) soll das Bündel frei
    von Flecken, Streifen und Hotspots sein. Auf einem Raster von \ang{1} quer und \ang{0.3} auf und ab
    werden die zweiten Differenzen zwischen Nachbarn in der unscharfen Ansicht genommen; ihr
    Effektivwert relativ zum Minimum von 50R (10\,100 cd) geht quadriert und gewichtet ins Ziel ein.
    Das Protokoll nennt ihn als \enquote{bumps}.
  \item[Die Suche.] Eine Zelle nach der anderen (die hellste zuerst): erst ihre Ablenkung über alle
    Möglichkeiten bei festem Zylinder, dann ihr Zylinder (Stärke und Richtung) bei fester Ablenkung, in
    Durchgängen über alle Zellen; nach dem ersten Durchgang nur noch Möglichkeiten nahe den
    bisherigen.
  \item[Entlang der Kanten.] \gui{Optimise: along the edges} (Standard an): wie die Riefen einer
    klassischen Scheibe streut und lenkt jedes Element längs der Kante seines eigenen Lichts -- des
    Asts seines Sektors. Ist die Sektortabelle leer, wird die Scheibe zuerst geschnitten: wo auf der
    ebenen Scheibe (in Schritten von \ang{1} Azimut) mehr als die Hälfte des Lichts über dem
    waagerechten Ast auf der Fahrerseite liegt, endet es am ansteigenden Ast -- der Keil vor dem
    Spiegelsektor, den die schräge Schirmkante beleuchtet (H4: von \ang{316} bis \ang{4}); der Rest
    endet am waagerechten Ast. Die Elemente des Keils bekommen um den Anstieg gedrehte Zylinder und
    Ablenkungen längs des Anstiegs, die übrigen Zylinder quer und jede Ablenkung. Ohne diese Bindung
    zieht die Suche das Licht des Keils zur Seite des Gegenverkehrs, wo R112 es zählt, und lässt den
    Anstieg dunkel.
  \item[Der Start.] \gui{Optimise: start with}: jede Zelle beginnt als Zylinderlinse mit \ang{6}
    Streuung quer (\gui{cylinder lenses (6\si{\degree} across)}, Standard), als ebenes Glas
    (\gui{plane glass}), oder beides wird versucht und das bessere behalten (\gui{both, the better
    kept}, doppelte Zeit). Die Suche endet nahe ihrem Anfang: aus ebenem Glas wird fast jede Zelle ein
    Prisma, aus Zylinderlinsen behalten viele ihre.
  \item[Runden.] \gui{Optimise: rounds} (Standard 2): die Scheibe wird gebaut und ihre Wendelbilder
    genommen; die halbe Differenz zur Summe der Muster wird zurückgeführt, die Zellen werden erneut
    optimiert, die beste Runde wird behalten.
  \item[Die Unschärfe der Lampe.] \gui{Optimise: the lamp's blur} = \gui{measured by ray tracing}
    (Standard): die Lampe, wie sie ist (Helix, kühlere Enden, Kolben), macht das Bündel weicher als
    die Wendelbilder -- bei der ebenen Scheibe einer H4 um $\sigma = \ang{0.5}$ quer und \ang{0.2}
    auf und ab. Das Plugin misst das zu Beginn per Raytracing an der ebenen Scheibe und beurteilt das
    Bündel danach in zwei Ansichten: jede Untergrenze in der dunkleren, jede Obergrenze in der
    helleren. Die behaltene Scheibe wird anschließend mit mindestens 20 Millionen Lichtwegen
    geraytract; was das Raytracing über die unscharfe Ansicht hinaus zeigt (ein schwacher Schleier),
    wird zurückgeführt und die Zellen werden noch zweimal optimiert. Behalten wird die beste nach dem
    Raytracing. \gui{none} schaltet das ab.
\end{description}

\subsection{Das Ergebnis}

\gui{Optimise lens} schreibt
\begin{itemize}
  \item die Elementtabelle neu: eine Zeile je Zelle (je Teil einer Zelle, die Sektoren schneiden),
  \item die Sektortabelle, wenn sie leer war (das Protokoll nennt die Sektoren),
\end{itemize}
und zeigt die Ergebnissätze \enquote{Low beam <Lampe> – optimised lens, filament images} und -- mit
gemessener Unschärfe -- \enquote{…, ray tracing}, je mit Bild, Prüfpunkten und Schnitten. Das Protokoll
nennt jede Runde (Werte der Ersatzrechnung und der Wendelbilder in cd, \enquote{bumps}), was behalten
wurde, ob R112 erfüllt ist, Lage und Schärfe der Grenze, Zahl der Elemente und Laufzeit. Die
Statusleiste: \enquote{lens optimised (… elements); filament images: R112 met; ray tracing: R112 met}.

Danach erzeugt \gui{Generate reflector} die optimierte Scheibe als Flächen.

\begin{hinweis}
Die Entwickler geben für einen typischen Fall an (H4, Rahmen $170 \times 100$ mm, Zellen 4 mm --
1063 Stück --, Brennpunkt 6,7 mm hinter der Wendel, Schirm mit 0,5, Unschärfe gemessen, 2 Runden,
Gleichmäßigkeit 1, Start mit Zylinderlinsen): etwa 9 Minuten und \qty{0.75}{GB} Speicher; alle
Anforderungen von R112 sowohl in den Wendelbildern als auch im Raytracing erfüllt. Mit Sektoren und
\gui{along the edges} (Standard) entstehen 1110 Elemente statt 1063 Zellen.
\end{hinweis}

\paragraph{Was die Einstellungen bewirken} (nach den Vergleichsrechnungen in der README des Plugins; \enquote{Standard}: mit Sektoren, entlang der Kanten, Band und Anstieg; \enquote{Band}: Mittel unter der Grenze, H $2 \dots 6^\circ$ ist der Anstieg):

\begin{center}
\small
\begin{tabular}{lrr}
\toprule
im Raytracing & Standard & frei, nur R112 \\
\midrule
Band, H $-10 \dots -2^\circ$ / $-2 \dots 0^\circ$ & 6700 / 14\,000 cd & 2500 / 7400 cd \\
Band, H $0 \dots 2^\circ$ / $2 \dots 6^\circ$ & 12\,400 / 3300 cd & 11\,800 / 3700 cd \\
75R / 50R / 50V & 12\,400 / 16\,900 / 13\,300 cd & 11\,000 / 19\,500 / 11\,900 cd \\
25L / Zone III & 2200 / 582 cd & 2600 / 594 cd \\
höchstens & 19\,000 cd & 20\,200 cd \\
Grenze & $-0{,}70^\circ$, $G = 0{,}28$ & $-0{,}70^\circ$, $G = 0{,}46$ \\
bumps (von 50R) & 3,8\,\% & 3,7\,\% \\
\bottomrule
\end{tabular}
\end{center}

Frei sammelt sich das Licht in einem Fleck 2 bis $3^\circ$ unter der Grenze; mit dem Band liegt es an der
Grenze und fällt nach unten ab. Aus ebenem Glas wird das Bündel unruhiger, ohne Gleichmäßigkeit sind
die Unebenheiten doppelt so hoch. Ohne Unschärfe erfüllen die Wendelbilder R112, das Raytracing findet
aber Licht in Zone III am Knick. Mit dem Brennpunkt nahe der Wendel (3 mm) sendet der glänzende Schirm
Licht über die Grenze, das keine Scheibe zurückholt (Zone III etwa 5500 cd).

% ---------------------------------------------------------------------------------------------
\section{Arbeitsabläufe Schritt für Schritt}\label{sec:al-ablaeufe}

\subsection{Ein erster Klarglasreflektor (H7)}\label{sec:al-ablauf-h7}

\begin{schritte}
  \item Öffnen Sie ein neues Dokument, das Dock \gui{Plugin} und wählen Sie \gui{Automotive
    Lighting}. Drücken Sie \gui{Reset parameters}: System Klarglas, Lampe H7, Rechtsverkehr,
    Rahmen $170 \times 100$ mm, Brennweite 25 mm, Standardzonen.
  \item Drücken Sie \gui{Generate reflector}. Nach einigen Sekunden bis Minuten (Fortschritt in der
    Statusleiste) erscheint der Spiegel in der Gruppe \enquote{Reflector H7}. Lesen Sie im \gui{Plugin
    log} Fläche, Tiefe, Abweichung der Wendelbilder von der Grenze, höchste Stufe und kleinsten
    Entformungswinkel.
  \item Drücken Sie \gui{Filament images}. Im Fenster \gui{Image} liegen alle Umrisse unter der
    hellen Linie der Grenze.
  \item Drücken Sie \gui{Light distribution}. Prüfen Sie im Fenster \gui{Table} die Spalte mit dem
    Urteil: \code{ok} oder \code{not met}. Im \gui{Plot} sehen Sie den Schnitt durch die Grenze und den
    Schnitt bei V $= -0{,}86^\circ$.
  \item Drücken Sie \gui{Light distribution (ray tracing)} und vergleichen Sie die beiden
    Ergebnissätze im Dock \gui{Results}.
\end{schritte}

\subsection{Den Reflektor an einen eigenen Rahmen anpassen}

\begin{schritte}
  \item Zeichnen Sie den Umriss des Spiegels in der Vorderansicht (Y-Z-Ebene) als geschlossene
    Kurve(n) -- relativ zur Wendelmitte, die Lampenachse bei Y = Z = 0 (bzw. bei \gui{Filament
    centre}).
  \item Wählen Sie die Kurven aus und drücken Sie bei \gui{Frame (curves)} auf \gui{Use selection}.
  \item Prüfen Sie \gui{Lamp hole radius} -- das Loch muss in den Rahmen passen und groß genug für
    die Verbindungsflächen sein.
  \item Drücken Sie \gui{Generate reflector}. Meldet das Plugin, dass die Kurven sich nicht schließen,
    verbinden Sie die Enden (Objektfang) und wählen sie neu.
  \item Ändern Sie danach eine Kurve, wählen Sie sie erneut mit \gui{Use selection}; eine gelöschte
    Kurve meldet \enquote{Frame: a chosen curve no longer exists – choose the frame again.}
\end{schritte}

\subsection{Zonen abstimmen}

\begin{schritte}
  \item Beginnen Sie mit \gui{Default zones}.
  \item Fehlt Licht am Hotspot (50R, 75R), verengen Sie die Streuung der Flanke mit dem waagerechten
    Ast (\gui{aim, right border} / \gui{aim, left border} näher zusammen) oder richten Sie sie weiter
    nach rechts.
  \item Ist das Vorfeld zu hell oder Zone I überschritten, setzen Sie bei den streuenden Zonen ein
    kleines \gui{below cut-off} nur, wo es gewollt ist, oder verringern Sie ihre Streuung.
  \item Meldet das Plugin \enquote{aim held back}, zielt eine Zone zu weit seitlich für ihren Platz
    am Spiegel: verringern Sie die Ziele oder verschieben Sie die Zonengrenzen.
  \item Prüfen Sie nach jeder Änderung mit \gui{Filament images} und \gui{Light distribution}.
  \item Für gleichmäßigere Streuung über lange Streifen setzen Sie \gui{Facet length} (etwa 20 bis
    30 mm) -- es entstehen mehr Verbindungsflächen.
\end{schritte}

\subsection{Ein H4-Scheinwerfer mit Streuscheibe}\label{sec:al-ablauf-h4}

\begin{schritte}
  \item Setzen Sie \gui{System} = \gui{patterned lens: paraboloid and lens}, \gui{Lamp} =
    \gui{H4 (low beam)}.
  \item Setzen Sie \gui{Mirror focus behind the filament centre} = 6.7 (Brennpunkt an der
    Fernlichtwendel) und z.\,B. \gui{Element width} = \gui{Element height} = 4.
  \item Prüfen Sie Lage und Dicke der Scheibe (\gui{Lens in front of the filament}, \gui{Thickness}).
  \item Drücken Sie \gui{Generate reflector}, dann \gui{Light distribution}: mit der Standardzeile
    der Elementtabelle (ein Kissen über die ganze Scheibe) sehen Sie die Grundverteilung.
  \item Drücken Sie \gui{Optimise lens}. Mit gemessener Unschärfe dauert das einige Minuten;
    \gui{Cancel} in der Statusleiste bricht ab.
  \item Sehen Sie sich die Ergebnissätze \enquote{… – optimised lens, filament images} und \enquote{…,
    ray tracing} an. Die Element- und Sektortabelle sind neu gefüllt.
  \item Drücken Sie \gui{Generate reflector}, um die optimierte Scheibe als Flächen zu erzeugen.
\end{schritte}

\begin{tipp}
Wollen Sie eine Optimierung mit anderer Einstellung wiederholen, leeren Sie die Sektortabelle, wenn
sie neu geschnitten werden soll -- \gui{Optimise lens} füllt sie nur, wenn sie leer ist.
\end{tipp}

\subsection{Eine Schlussleuchte prüfen (P21W)}

\begin{schritte}
  \item \gui{System} = \gui{patterned lens: paraboloid and lens}, \gui{Lamp} = \gui{P21W}
    (Querwendel).
  \item \gui{Test points} = \gui{stop lamp (S1)}.
  \item Setzen Sie \gui{Transmittance} auf den Wert der roten Scheibe.
  \item Stellen Sie die Elemente so ein, dass das Licht etwa $\pm 20^\circ$ horizontal und
    $\pm 10^\circ$ vertikal verteilt wird (z.\,B. ein Kissen mit großer Streuung).
  \item Drücken Sie \gui{Generate reflector} und \gui{Light distribution}. Die Tabelle \enquote{R148
    standard light distribution, stop lamp (S1)} zeigt jeden Punkt in Prozent des Minimums auf der
    Achse und das Maximum.
\end{schritte}

\subsection{Die Lampe ansehen und das Fernlicht der H4 prüfen}

\begin{schritte}
  \item Wählen Sie die H4 und drücken Sie \gui{Generate lamp}.
  \item Blenden Sie im Modellbaum \enquote{Bulb glass} und den Wendelzylinder aus; zoomen Sie auf die
    Wendel. Die Attribute \code{al.luminance} der Drahtstücke zeigen die kühleren Enden.
  \item Setzen Sie \gui{H4: lit coil} = \gui{driving beam} und drücken Sie \gui{Light distribution
    (ray tracing)}: Ergebnissatz \enquote{Driving beam H4 – ray tracing} mit den Fernlichtpunkten.
  \item Setzen Sie \gui{H4: lit coil} wieder auf \gui{passing beam}.
\end{schritte}

\subsection{Ein Nebelscheinwerfer (H3)}

\begin{schritte}
  \item \gui{Lamp} = \gui{H3}, \gui{Rise} = 0.
  \item Ersetzen Sie die Zonentabelle durch eine Zone (\gui{from azimuth} 0, \code{horizontal},
    Ziele z.\,B. 20 und $-20$) oder eine obere und eine untere.
  \item \gui{Generate reflector}, \gui{Filament images}, \gui{Light distribution}.
\end{schritte}

% ---------------------------------------------------------------------------------------------
\section{Typische Fehler und Meldungen}\label{sec:al-fehler}

\begin{referenz}{>{\raggedright\arraybackslash}p{6.4cm}>{\raggedright\arraybackslash}X}
\toprule
Meldung (gekürzt) & Abhilfe \\
\midrule
\endhead
\code{The zone table is empty …} & Zone anlegen oder \gui{Default zones}. \\
\code{Two zones start at the same azimuth.} & Azimute der Zonen unterschiedlich wählen. \\
\code{the frame curves do not join to one outline (gap of … mm)} & Kurvenenden verbinden. \\
\code{the frame is not closed …} / \code{the frame has no area} & geschlossenen Umriss wählen. \\
\code{no reflector area: the zones do not cover the frame outside the lamp hole} & Zonen und Rahmen
  prüfen. \\
\code{a cell of the reflector would have a round hole …} & schmalere Streifen als das Lampenloch oder
  weniger seitliche Ziele. \\
\code{a strip border does not reach this section …} & Zone schmaler machen oder ihr Ziel verringern. \\
\code{this part of the mirror cannot throw its light where the zone is aimed …} & Zone weniger weit
  seitlich zielen lassen. \\
\code{the strips of this zone run along its start section …} & Ziel passt nicht zur Lage der Zone. \\
\code{a facet becomes too steep or folds …} & Streuung, Streifenbreite oder Zonenausdehnung
  verringern. \\
\code{the strips of a facet cross …} & Ziele der Zone andersherum laufen lassen oder Streuung
  verringern. \\
\code{a facet for a lamp with its filament across the axis cannot keep to one side …} & schmalere
  Streifen oder weniger weite Ziele. \\
\code{the strips are too narrow (… mm) for a connecting face over a step of … mm} & breitere
  Streifen, kleinere Schräge oder kleinerer Verrundungsradius. \\
\code{the connecting faces need more room than the facets can give …} & Schräge oder
  Verrundungsradius verringern. \\
\code{the lamp hole (radius … mm) is too small …: at least … mm} & \gui{Lamp hole radius}
  vergrößern. \\
\code{the lens (at x = … mm on the axis) lies inside the mirror …} & Scheibe weiter nach vorn oder
  weniger neigen. \\
\code{the lens' radius (… mm) is smaller than the lens is wide} & \gui{Dome radius} vergrößern. \\
\code{the lens element at y = …, z = … mm cannot bend the light that far …} & Ablenkung oder
  Streuung verringern. \\
\code{the lens elements are … mm deep, more than the lens is thick …} & dickere Scheibe, kleinere
  Zellen, weniger Ablenkung. \\
\code{walls with a draft angle and lens sectors together are not supported yet …} & Schräge 0 oder
  Sektoren leeren. \\
\code{the lens elements are too small for walls with this draft angle and fillet radius …} & größere
  Zellen, kleinere Schräge oder Verrundung. \\
\code{the lens is too thin for …} (\gui{Optimise lens}) & dickere Scheibe. \\
\code{Strips belong to the facetted reflector …} & \gui{Show strips} nur beim Klarglasreflektor. \\
\code{the C++ kernel (package pathtracer 0.3.0 or later) is not installed …} & \gui{Kernel} auf
  \code{automatic} oder \code{pathtracer} installieren. \\
\code{the ray tracing needs the Path Tracer plugin next to this one …} & Path-Tracer-Plugin
  installieren (NumPy-Kern). \\
\code{the coil (or shield) does not fit into a bulb of … mm …} & \gui{Bulb: outer diameter} oder
  \gui{Bulb: wall} anpassen (0: typisch). \\
\code{the bulb's wall … must be thinner than a quarter of its diameter} & dünnere Wand. \\
\bottomrule
\end{referenz}

% ---------------------------------------------------------------------------------------------
\section{Grenzen}\label{sec:al-grenzen}

\begin{itemize}
  \item \textbf{Photometrie} (\gui{Filament images}, \gui{Light distribution}): die Wendel als
    gleichmäßig leuchtender Zylinder, eine Reflexion an Facetten und Verbindungsflächen, kein
    Kolbenglas, keine Abschlussscheibe, kein Gehäuse; von den inneren Teilen der Lampe nur Schirm und
    Kappe der H4 (nicht die Fernlichtwendel und Halter, nicht das Licht, das die Schirminnenseite
    reflektiert). Das Direktlicht auf einem Raster von \ang{1}, von der Wendelmitte aus. Die Grenze ist
    schärfer als bei einer echten Lampe. Mit \gui{Shadows of the steps} wird keine dritte Reflexion
    verfolgt; der Spiegel wird als Höhenkarte von 0,25 mm verglichen.
  \item \textbf{Raytracing}: typische Kolben, hinten offen (kein Quetschfuß), keine Halter oder
    Zuleitungen der Wendel; Endtemperatur angenommen, die Windungen wärmen sich nicht gegenseitig; die
    Drähte absorbieren, was sie trifft (Wolfram reflektiert etwa die Hälfte); der Schirm ist ein idealer
    Spiegel; die Scheibe hat keinen Rand.
  \item \textbf{Streuscheibe}: ein an einer Wand gespiegelter Strahl trifft keine zweite Wand (sonst
    verloren); Reflexionen an den beiden Scheibenseiten fehlen in der Photometrie. Die Kuppel ist eine
    Kugel um einen Punkt der Lampenachse; die Elemente stehen in X-Richtung auf der Scheibenform.
  \item Das Lampenloch muss groß genug für die Bänder zwischen den Zonen sein; die Streifen müssen
    breiter sein als ihre Verbindungsflächen.
  \item Der Rahmen sollte konvex oder nahezu konvex sein; die Lampenachse ist $+X$; die Flächen sind
    einzelne getrimmte Flächen, nicht zu einer Schale vernäht.
\end{itemize}
